% 2020 天津卷 第 4 题：频率分布直方图
% 来源 PDF：普通高考/2020/2020天津.pdf 第 1 页（A3 横排），印面 1131，第 4 题题干下方
% 几何按原页矢量数据换算：9 个等宽组 [5.31,5.33) … [5.47,5.49]，
%   矩形高度（频率/组距）依次为 1.25、3.75、2.50、7.50、10.00、8.75、6.25、5.00、5.00，
%   组距 0.02，故各矩形面积（频率）依次为 0.025、0.075、0.05、0.15、0.20、0.175、0.125、0.10、0.10，
%   合计 1.00，与 80 个样本一致。
% 纵轴虚线只标 1.25 的整数倍刻度，自 1.25 起至 10.00（原页即 8 条虚线）；
% 横轴自 5.31 起每 0.02 标一组端点，末尾 5.49；原点左侧有断裂符号，0 标在断裂符号之外。
\begin{tikzpicture}[font=\small,
  x=0.60cm, y=0.26cm, line cap=round, line join=round]
  % 坐标轴（左端留出断裂符号与 0 的位置）
  \draw[->] (-1.35,0)--(10.35,0) node[right,inner sep=4pt] {直径/mm};
  \draw[->] (0,0)--(0,11.6) node[above,inner sep=4pt] {$\dfrac{\text{频率}}{\text{组距}}$};
  % 断裂符号
  \draw (-0.78,0.10)--(-0.36,-0.34) (-0.78,-0.34)--(-0.36,0.10);
  \node[below,inner sep=4pt] at (-1.18,0) {$0$};
  % 纵向虚线刻度 1.25—10.00
  \foreach \h/\lab in {1.25/1.25,2.50/2.50,3.75/3.75,5.00/5.00,
                       6.25/6.25,7.50/7.50,8.75/8.75,10.00/10.00}{
    \draw[dashed,line width=0.4pt] (0,\h)--(9,\h);
    \draw (-0.14,\h)--(0,\h);
    \node[left,inner sep=2pt,font=\scriptsize] at (-0.12,\h) {\lab};
  }
  % 9 个矩形
  \foreach \i/\h in {0/1.25,1/3.75,2/2.50,3/7.50,4/10.00,5/8.75,6/6.25,7/5.00,8/5.00}{
    \draw[line width=0.6pt] (\i,0) rectangle (\i+1,\h);
  }
  % 横轴刻度与标签
  \foreach \i/\lab in {0/5.31,1/5.33,2/5.35,3/5.37,4/5.39,5/5.41,
                       6/5.43,7/5.45,8/5.47,9/5.49}{
    \draw (\i,0)--(\i,-0.28);
    \node[below,inner sep=2pt,font=\scriptsize] at (\i,0) {\lab};
  }
\end{tikzpicture}
